This introduction defines generative diffusion, anchors its empirical impact, and previews the survey across §2--§14. Generative diffusion models (GDMs) learn a data distribution \(p_d\)ata(x) by reversing a fixed Markovian noising process. A forward chain perturbs a clean datum \(x_0\) with Gaussian noise on a schedule β\_1, \ldots, β\_T until \(x_T\) is approximately N(0, I), and a neural network learns to invert that corruption step by step. The recipe is summarized as ``creating noise from data is easy; creating data from noise is generative modeling'' (Song et al., ICLR 2021). It now dominates high-fidelity controllable generation across images, video, audio, 3D, motion, molecules, proteins, and decision making, and by 2026 underlies every major commercial text-to-image system. The Stable Diffusion family progresses from a latent U-Net at 512² in 1.5 (2022), through OpenCLIP-H/14 conditioning in 2.1 (2022) and a 2.6B U-Net at 1024² in SDXL (2023), to MMDiT with rectified flow in SD3 (2024). Closed systems include DALL-E 2 (2022, unCLIP cascade), DALL-E 3 (2023, synthetic-caption training), Imagen (2022, T5-XXL cascade), and Midjourney v6 (2023, proprietary photorealism). Open variants include Kandinsky 2.2 (2022, image-prior diffusion), PixArt-α (2024, 600M DiT), Flux.1 (2024, 12B MMDiT), and Hunyuan-DiT (2024, bilingual DiT). In video, Sora (2024) introduced spatiotemporal-patch DiT, Stable Video Diffusion (2023) shipped a 1.5B 3D U-Net, and W.A.L.T (2023) demonstrated causal-VAE vi\ldots{}
